\subsection{DIHEDRAL GROUPS}

DEFINITION 2. A dihedral group is a group generated by two distinct elements of order two.

Example. Let M be the multiplicative group \(\{1,-1\}\) , and let m be an integer \(\geqslant2\) (resp. \(m=\infty\) ). Then M acts on the group Z/mZ (resp. on Z) by \((-1).x=-x\) , and the corresponding semi-direct product of M by Z/mZ (resp. of M by Z) is denoted by \(D_{m}\) . The elements of \(D_{m}\) are thus the pairs \((\varepsilon,x)\) with \(\varepsilon=\pm1\) and \(x\in Z/mZ\) (resp. \(x\in Z\) ); the group law on \(D_{m}\) is given by the formula

\[
(\varepsilon , x). (\varepsilon^ {\prime}, x ^ {\prime}) = (\varepsilon \varepsilon^ {\prime}, \varepsilon^ {\prime} x + x ^ {\prime}).\tag{6}
\]

We denote by \(\iota\) the class of 1 modulo \(m\) (resp. \(\iota = 1\) ) and set

\[
\rho = (- 1, 0), \quad \rho^ {\prime} = (- 1, \iota), \quad \pi = (1, \iota).\tag{7}
\]

Then \(\rho^2 = \rho'^2 = 1\) and \(\pi = \rho \rho'\) . The formulas

\[
\pi^ {n} = (1, n \iota), \quad \rho \pi^ {n} = (- 1, n \iota)\tag{8}
\]

show that \(D_{m}\) is a dihedral group generated by \(\{\rho,\rho'\}\) .

PROPOSITION 2. Assume that \(S\) consists of two distinct elements \(s\) and \(s'\) of order 2.

\begin{enumerate}
\def\labelenumi{(\roman{enumi})}
\item
  The subgroup P of W generated by \(p = ss'\) is normal, and W is the semi-direct product of the subgroup \(T = \{1, s\}\) and P. Moreover, \((W : P) = 2\) .
\item
  Let \(m\) be the order (finite or infinite) of \(p\) . Then \(m \geqslant 2\) and \(W\) is of order \(2m\) . There is a unique isomorphism \(\varphi\) from \(D_m\) to \(W\) such that \(\varphi(\rho) = s\) and \(\varphi(\rho') = s'\) .
\item
  We have \(sps^{-1} = sss's = s's = p^{-1}\) , and hence
\end{enumerate}

\[
s p ^ {n} s ^ {- 1} = p ^ {- n}\tag{9}
\]

for every integer \(n\) . Since \(W\) is generated by \(\{s, s'\}\) , and hence by \(\{s, p\}\) , the subgroup \(P\) is normal. It follows that \(TP\) is a subgroup of \(W\) , and since \(TP\) contains \(s\) and \(s' = sp\) , we have \(W = TP = P \cup sP\) . To prove (i), it is therefore enough to show that \(W \neq P\) . If \(W = P\) , the group \(W\) would be commutative, so \(p^2 = s^2 s'^2 = 1\) . But then the only elements of \(W = P\) would be 1 and \(p\) , contradicting the hypothesis that \(W\) contains at least three elements, namely 1, \(s\) and \(s'\) .

\begin{enumerate}
\def\labelenumi{(\roman{enumi})}
\setcounter{enumi}{1}
\tightlist
\item
  Since \(s \neq s'\) , we have \(p \neq 1\) and so \(m \geqslant 2\) . Since P is of order m and (W : P) = 2, the order of W is 2m. If m is finite (resp. infinite), there exists an isomorphism \(\varphi'\) from Z/mZ (resp. Z) to P taking \(\pi\) to p.~Moreover, there exists an isomorphism \(\varphi''\) from M = \{1, -1\} to T taking -1 to s. The group W is the semi-direct product of T and P. In view of the formulas (9) and \(\rho\pi^{n}\rho^{-1} = \pi^{-n}\) , \(\varphi'\) and \(\varphi''\) induce an isomorphism \(\varphi\) from D\_\{m\} to W such that \(\varphi(\rho) = s\) and \(\varphi(\pi) = p\) , and hence \(\varphi(\rho') = s'\) . The uniqueness of \(\varphi\) follows from the fact that D\_\{m\} is generated by \(\{\rho, \rho'\}\) .
\end{enumerate}

Remark. Consider a dihedral group W of order 2m generated by two distinct elements s and \(s'\) of order 2. Denote by \(s_{q}\) (resp. \(s'_{q}\) ) the sequence of length q whose odd (resp. even) numbered terms are equal to s and whose even (resp. odd) numbered terms are equal to \(s'\) , and let \(w_{q}\) (resp. \(w'_{q}\) ) be the product of the sequence \(s_{q}\) (resp. \(s'_{q}\) ). We have

\[
\begin{array}{r l} w _ {2 k} = (s s ^ {\prime}) ^ {k}, & w _ {2 k + 1} = (s s ^ {\prime}) ^ {k} s, \\ w _ {2 k} ^ {\prime} = (s ^ {\prime} s) ^ {k} = (s s ^ {\prime}) ^ {- k}, & w _ {2 k + 1} ^ {\prime} = (s ^ {\prime} s) ^ {k} s ^ {\prime} = (s s ^ {\prime}) ^ {- k - 1} s. \end{array}
\]

If \(\mathbf{s} = (s_1, \ldots, s_q)\) is a reduced decomposition (with respect to \(\{s, s'\}\) ) of an element \(w\) of W, then clearly \(s_i \neq s_{i+1}\) for \(1 \leqslant i \leqslant q-1\) . Hence, \(\mathbf{s} = \mathbf{s}_q\) or \(\mathbf{s} = \mathbf{s}_q'\) .

If \(m = \infty\) , the elements \((ss')^n\) and \((ss')^ns\) for \(n \in \mathbf{Z}\) are distinct. Hence, the elements \(w_q\) ( \(q \geqslant 0\) ) and \(w_q'\) ( \(q > 0\) ) are distinct, and if \(\mathbf{s}\) is a reduced decomposition of \(w_q\) (resp. \(w_q'\) ) we necessarily have \(\mathbf{s} = \mathbf{s}_q\) (resp. \(\mathbf{s} = \mathbf{s}_q'\) ). It follows from this that \(l(w_q) = l(w_q') = q\) and that the set of reduced decompositions of the elements of \(\mathbf{W}\) consists of the \(\mathbf{s}_q\) and the \(\mathbf{s}_q'\) . Moreover, every element of \(\mathbf{W}\) has a unique reduced decomposition.

Suppose now that \(m\) is finite. If \(q \geqslant 2m\) , we have \(w_{q} = w_{q - 2m}\) and \(w_{q}' = w_{q - 2m}'\) ; if \(m \leqslant q \leqslant 2m\) , we have \(w_{q} = w_{2m - q}'\) , \(w_{q}' = w_{2m - q}\) . Hence, neither \(\mathbf{s}_q\) nor \(\mathbf{s}_q'\) are reduced decompositions if \(q > m\) . It follows that each of the \(2m\) elements of \(W\) is one of the \(2m\) elements \(w_0 = w_0'\) , \(w_q\) and \(w_q'\) for \(1 \leqslant q \leqslant m - 1\) , and \(w_m = w_m'\) . These \(2m\) elements are thus distinct and it follows as above that \(l(w_q) = l(w_q') = q\) for \(q \leqslant m\) and that the set of reduced decompositions of elements of \(W\) consists of the \(s_q\) and the \(s_q'\) for \(0 \leqslant q \leqslant m\) . Every element of \(W\) except \(w_m\) has a unique reduced decomposition; \(w_m\) has two.

